% Three legal execution orders of the six blocks of a grid: in index order,
% in an arbitrary order, and two at a time.
% Redraws the outline figure m3l1_execution_order_blocks.
% Usage: % Three legal execution orders of the six blocks of a grid: in index order,
% in an arbitrary order, and two at a time.
% Redraws the outline figure m3l1_execution_order_blocks.
% Usage: % Three legal execution orders of the six blocks of a grid: in index order,
% in an arbitrary order, and two at a time.
% Redraws the outline figure m3l1_execution_order_blocks.
% Usage: % Three legal execution orders of the six blocks of a grid: in index order,
% in an arbitrary order, and two at a time.
% Redraws the outline figure m3l1_execution_order_blocks.
% Usage: \input{figures/l03-block-order}   (body width, 120 mm)
\begin{tikzpicture}[x=1mm, y=1mm, font=\scriptsize\sffamily,
                    >={Stealth[length=1.4mm]},
                    blk/.style={draw, fill=ln-box, minimum width=13mm,
                                minimum height=4.5mm, inner sep=1pt}]
  % the grid
  \draw[ln-accent, thick, rounded corners=1pt] (0,14) rectangle (22,30);
  \node[anchor=south, text=ln-accent] at (11,30.3) {\iflangja{グリッド}{grid}};
  \foreach \bx in {0,1,2} \foreach \by in {0,1} {
    \draw[fill=ln-box] (1.5+\bx*6.6,16+\by*6.5) rectangle ++(6,5.5);
    \node[font=\scriptsize\sffamily] at (4.5+\bx*6.6,18.7+\by*6.5) {\bx,\by};
  }
  \draw[->, ln-muted] (23,22) -- (27,22);
  % time axis
  \draw[->, ln-muted] (30,-2) -- (118,-2) node[below left, text=ln-muted]
    {\iflangja{時間}{time}};
  % case 1
  \node[anchor=east, text=ln-muted, align=right] at (29,33)
    {\iflangja{順序 1}{case 1}};
  \foreach \k/\t in {0/{0,0}, 1/{1,0}, 2/{2,0}, 3/{0,1}, 4/{1,1}, 5/{2,1}} {
    \node[blk, anchor=west] at (30+\k*14.5,33) {\t};
  }
  % case 2
  \node[anchor=east, text=ln-muted, align=right] at (29,25)
    {\iflangja{順序 2}{case 2}};
  \foreach \k/\t in {0/{2,1}, 1/{1,1}, 2/{1,0}, 3/{0,0}, 4/{2,0}, 5/{0,1}} {
    \node[blk, anchor=west] at (30+\k*14.5,25) {\t};
  }
  % case 3: three at a time
  \node[anchor=east, text=ln-muted, align=right] at (29,10)
    {\iflangja{順序 3}{case 3}};
  \foreach \k/\t in {0/{0,0}, 1/{1,0}, 2/{2,0}} {
    \node[blk, anchor=west, minimum width=40mm] at (30,16-\k*5.5) {\t};
  }
  \foreach \k/\t in {0/{0,1}, 1/{1,1}, 2/{2,1}} {
    \node[blk, anchor=west, minimum width=40mm] at (74,16-\k*5.5) {\t};
  }
\end{tikzpicture}
   (body width, 120 mm)
\begin{tikzpicture}[x=1mm, y=1mm, font=\scriptsize\sffamily,
                    >={Stealth[length=1.4mm]},
                    blk/.style={draw, fill=ln-box, minimum width=13mm,
                                minimum height=4.5mm, inner sep=1pt}]
  % the grid
  \draw[ln-accent, thick, rounded corners=1pt] (0,14) rectangle (22,30);
  \node[anchor=south, text=ln-accent] at (11,30.3) {\iflangja{グリッド}{grid}};
  \foreach \bx in {0,1,2} \foreach \by in {0,1} {
    \draw[fill=ln-box] (1.5+\bx*6.6,16+\by*6.5) rectangle ++(6,5.5);
    \node[font=\scriptsize\sffamily] at (4.5+\bx*6.6,18.7+\by*6.5) {\bx,\by};
  }
  \draw[->, ln-muted] (23,22) -- (27,22);
  % time axis
  \draw[->, ln-muted] (30,-2) -- (118,-2) node[below left, text=ln-muted]
    {\iflangja{時間}{time}};
  % case 1
  \node[anchor=east, text=ln-muted, align=right] at (29,33)
    {\iflangja{順序 1}{case 1}};
  \foreach \k/\t in {0/{0,0}, 1/{1,0}, 2/{2,0}, 3/{0,1}, 4/{1,1}, 5/{2,1}} {
    \node[blk, anchor=west] at (30+\k*14.5,33) {\t};
  }
  % case 2
  \node[anchor=east, text=ln-muted, align=right] at (29,25)
    {\iflangja{順序 2}{case 2}};
  \foreach \k/\t in {0/{2,1}, 1/{1,1}, 2/{1,0}, 3/{0,0}, 4/{2,0}, 5/{0,1}} {
    \node[blk, anchor=west] at (30+\k*14.5,25) {\t};
  }
  % case 3: three at a time
  \node[anchor=east, text=ln-muted, align=right] at (29,10)
    {\iflangja{順序 3}{case 3}};
  \foreach \k/\t in {0/{0,0}, 1/{1,0}, 2/{2,0}} {
    \node[blk, anchor=west, minimum width=40mm] at (30,16-\k*5.5) {\t};
  }
  \foreach \k/\t in {0/{0,1}, 1/{1,1}, 2/{2,1}} {
    \node[blk, anchor=west, minimum width=40mm] at (74,16-\k*5.5) {\t};
  }
\end{tikzpicture}
   (body width, 120 mm)
\begin{tikzpicture}[x=1mm, y=1mm, font=\scriptsize\sffamily,
                    >={Stealth[length=1.4mm]},
                    blk/.style={draw, fill=ln-box, minimum width=13mm,
                                minimum height=4.5mm, inner sep=1pt}]
  % the grid
  \draw[ln-accent, thick, rounded corners=1pt] (0,14) rectangle (22,30);
  \node[anchor=south, text=ln-accent] at (11,30.3) {\iflangja{グリッド}{grid}};
  \foreach \bx in {0,1,2} \foreach \by in {0,1} {
    \draw[fill=ln-box] (1.5+\bx*6.6,16+\by*6.5) rectangle ++(6,5.5);
    \node[font=\scriptsize\sffamily] at (4.5+\bx*6.6,18.7+\by*6.5) {\bx,\by};
  }
  \draw[->, ln-muted] (23,22) -- (27,22);
  % time axis
  \draw[->, ln-muted] (30,-2) -- (118,-2) node[below left, text=ln-muted]
    {\iflangja{時間}{time}};
  % case 1
  \node[anchor=east, text=ln-muted, align=right] at (29,33)
    {\iflangja{順序 1}{case 1}};
  \foreach \k/\t in {0/{0,0}, 1/{1,0}, 2/{2,0}, 3/{0,1}, 4/{1,1}, 5/{2,1}} {
    \node[blk, anchor=west] at (30+\k*14.5,33) {\t};
  }
  % case 2
  \node[anchor=east, text=ln-muted, align=right] at (29,25)
    {\iflangja{順序 2}{case 2}};
  \foreach \k/\t in {0/{2,1}, 1/{1,1}, 2/{1,0}, 3/{0,0}, 4/{2,0}, 5/{0,1}} {
    \node[blk, anchor=west] at (30+\k*14.5,25) {\t};
  }
  % case 3: three at a time
  \node[anchor=east, text=ln-muted, align=right] at (29,10)
    {\iflangja{順序 3}{case 3}};
  \foreach \k/\t in {0/{0,0}, 1/{1,0}, 2/{2,0}} {
    \node[blk, anchor=west, minimum width=40mm] at (30,16-\k*5.5) {\t};
  }
  \foreach \k/\t in {0/{0,1}, 1/{1,1}, 2/{2,1}} {
    \node[blk, anchor=west, minimum width=40mm] at (74,16-\k*5.5) {\t};
  }
\end{tikzpicture}
   (body width, 120 mm)
\begin{tikzpicture}[x=1mm, y=1mm, font=\scriptsize\sffamily,
                    >={Stealth[length=1.4mm]},
                    blk/.style={draw, fill=ln-box, minimum width=13mm,
                                minimum height=4.5mm, inner sep=1pt}]
  % the grid
  \draw[ln-accent, thick, rounded corners=1pt] (0,14) rectangle (22,30);
  \node[anchor=south, text=ln-accent] at (11,30.3) {\iflangja{グリッド}{grid}};
  \foreach \bx in {0,1,2} \foreach \by in {0,1} {
    \draw[fill=ln-box] (1.5+\bx*6.6,16+\by*6.5) rectangle ++(6,5.5);
    \node[font=\scriptsize\sffamily] at (4.5+\bx*6.6,18.7+\by*6.5) {\bx,\by};
  }
  \draw[->, ln-muted] (23,22) -- (27,22);
  % time axis
  \draw[->, ln-muted] (30,-2) -- (118,-2) node[below left, text=ln-muted]
    {\iflangja{時間}{time}};
  % case 1
  \node[anchor=east, text=ln-muted, align=right] at (29,33)
    {\iflangja{順序 1}{case 1}};
  \foreach \k/\t in {0/{0,0}, 1/{1,0}, 2/{2,0}, 3/{0,1}, 4/{1,1}, 5/{2,1}} {
    \node[blk, anchor=west] at (30+\k*14.5,33) {\t};
  }
  % case 2
  \node[anchor=east, text=ln-muted, align=right] at (29,25)
    {\iflangja{順序 2}{case 2}};
  \foreach \k/\t in {0/{2,1}, 1/{1,1}, 2/{1,0}, 3/{0,0}, 4/{2,0}, 5/{0,1}} {
    \node[blk, anchor=west] at (30+\k*14.5,25) {\t};
  }
  % case 3: three at a time
  \node[anchor=east, text=ln-muted, align=right] at (29,10)
    {\iflangja{順序 3}{case 3}};
  \foreach \k/\t in {0/{0,0}, 1/{1,0}, 2/{2,0}} {
    \node[blk, anchor=west, minimum width=40mm] at (30,16-\k*5.5) {\t};
  }
  \foreach \k/\t in {0/{0,1}, 1/{1,1}, 2/{2,1}} {
    \node[blk, anchor=west, minimum width=40mm] at (74,16-\k*5.5) {\t};
  }
\end{tikzpicture}
